% The Figure structure intentionally includes the non-floating caption.
\tagstructbegin{tag=Figure,alttext={训练与验证损失曲线不同步，提前停止选择验证误差较低时的参数}}
\begin{center}
\begin{tikzpicture}
\begin{axis}[
 width=112mm,height=46mm,
 xmin=.8,xmax=8.3,ymin=.2,ymax=1.28,
 axis lines=left,xlabel={训练轮次（Epoch）},ylabel={损失},
 xtick={1,2,3,4,5,6,7,8},
 tick label style={font=\footnotesize},
 label style={font=\small},
 legend style={at={(.98,.97)},anchor=north east,font=\footnotesize,
 fill=white,draw=linegray},
]
\addplot[blue,very thick,mark=*] coordinates {
(1,1.17) (2,.89) (3,.71) (4,.55) (5,.44) (6,.35) (7,.29) (8,.25)};
\addlegendentry{训练损失}
\addplot[orange,very thick,mark=square*] coordinates {
(1,1.22) (2,.96) (3,.80) (4,.74) (5,.78) (6,.84) (7,.94) (8,1.06)};
\addlegendentry{验证损失}
\draw[dashed,teal,thick] (axis cs:4,.2)--(axis cs:4,1.2);
\node[teal,anchor=south west,font=\footnotesize] at (axis cs:4.12,.23) {最佳验证轮次};
\end{axis}
\end{tikzpicture}

{\footnotesize 人工示意数据：训练损失持续下降，
验证损失在第 4 轮后反而上升；选检查点应看验证集。}
\end{center}
\begin{center}\captionof{figure}{训练损失与验证损失走势及最优验证轮次}\label{fig:early-stopping}\end{center}
\tagstructend
